% Copyright 2026 The terCAD team. All rights reserved.
% Use of this content is governed by a CC BY-NC-ND 4.0 license that can be found in the LICENSE file.

\subsection{Why Budgeting Matters}
\markboth{Budgeting}{Why Budgeting Matters}

Budget. For many people, the word itself suggests restriction: income always seems finite,
while the list of things we need (or simply want) keeps growing. However, a budget is not
about saying no to everything. It's about deciding in advance how to spend your money, rather
than discovering at the end of the month where it went.

Without a plan, spending tends to creep up: a subscription here, a dinner out there, and a
few unplanned purchases can easily cancel out any progress towards important goals. A budget
gives every income a purpose, turning vague intentions ('I should save more') into measurable
numbers that can be tracked day by day.

When building a budget, it is important to bear in mind the two finite resources of money and
time. Everything costs one or the other, and the two are closely connected. For example,
cooking at home saves money but takes time, while owning a car saves time but costs money. A
good budget acknowledges these trade-offs and makes them visible, enabling you to make
conscious decisions.

Finally, a budget is not a static document. It is not a one-off exercise, but rather a living
plan that is reviewed and adjusted in response to life's changes: a new job, a move, having a
child or paying off a loan. A budget designed for a student will not suit a family of four,
so there is no reason to keep it unchanged when circumstances change.


\subsection{Choosing a Strategy}
\markboth{Budgeting}{Choosing a Strategy}

A single magic formula is not known, but several well-known approaches are recognized. All of
these can be implemented in Fingrom using the tools described later in this section.

\subsubsection{The 50/30/20 Rule}

This rule, popularised by Elizabeth Warren and Amelia Warren Tyagi \cite{Warr05}, splits the
after-tax income into three parts:

\begin{itemize}
  \item 50\% -- needs: housing, utilities, groceries, transportation, insurance, minimum loan
  payments;
  \item 30\% -- wants: dining out, entertainment, hobbies, travel;
  \item 20\% -- savings and extra debt repayment.
\end{itemize}

\noindent In Fingrom, this maps directly to 'relative' limits (see \cref{budget:limit}): a
category with a limit of 0.5 always receives half of the income received in the current
period, regardless of its size.


\subsubsection{Zero-Based Budgeting}

This approach assumes that everything starts from zero, and that every unit of income must be
assigned to a category until nothing remains unallocated: \emph{Income} $-$ \emph{Expenses}
$-$ \emph{Savings} $= 0$. This method is particularly useful when starting from scratch or
after a significant change, as it prevents habits from carrying over. In Fingrom, it is
implemented with fixed limits whose sum equals the expected monthly income (\cref{img:fixed-income}).

\img{img/ch01-s01-budget__fixed-income}{Fixed Income Budgeting}{img:fixed-income}


\subsubsection{Incremental Budgeting}

Incremental budgeting uses the previous period as a basis and adjusts it by a certain
percentage (an \q{$x + z\%$} view). This approach is ineffective on the first day, but
becomes useful once a few months of data have been recorded. The \textbf{Summary} tab of each
category and the \textbf{Metrics} page (see \cref{budget:metrics}) provide the exact numbers
needed, showing how much was actually spent in each previous month.


\subsubsection{Activity-Based Budgeting}

This approach works in reverse: starting from a desired outcome (such as a holiday, a new
laptop or an emergency fund), it calculates backwards to determine what is needed to achieve
it. In Fingrom, the desired outcome is defined as a \textbf{Goal}, while the budget
categories ensure that there is enough room in the remaining expenses to reach it.


\subsubsection{Envelope Method}

The classic approach to budgeting, where cash is divided into physical envelopes, one per
spending category, is now available digitally. Once an envelope is empty, spending in that
category stops until the next period. Each Fingrom budget category behaves exactly like such
an envelope: the remaining amount is always displayed next to the category, and exceeding the
limit is clearly highlighted (\cref{img:budget-list}).

\img{img/ch01-s01-budget__budget-list}{Budget List}{img:budget-list}


\subsubsection{Characteristics of a Good Budget}

Regardless of the chosen strategy, the budget should:

\begin{enumerate}
  \item \textbf{Realistic.} One of the fastest ways to abandon a budget is to create an
  unrealistic one. Start from actual spending, not an idealized version of it.
  \item \textbf{Flexible.} It should be easy to adjust when income or priorities change.
  Percentages (relative limits) scale up and down automatically, whereas fixed amounts need to be revised manually.
  \item \textbf{Easily understood.} A dozen well-named categories are more useful than a
  hundred precise ones. If it takes too long to choose a category for a bill, the structure is too complex.
  \item \textbf{Mindful of irregularities.} Not all expenses or income are monthly.
  Insurance, taxes, holidays and annual subscriptions should have their own place, either as
  yearly categories or as monthly coefficients (see \cref{budget:coef}).
\end{enumerate}

\noindent The resources and the way to use them are unique to each person. Blindly copying
someone else's recommendation rarely works, since it does not take personal circumstances or
risk tolerance into account. Use the above strategies as a starting point and allow recorded
history to inform subsequent adjustments.


\subsection{Budget Categories}
\markboth{Budgeting}{Budget Categories}

A budget category is a specific area to which money is allocated in order to track expenses,
such as groceries, rent, transport and leisure. Every bill recorded in Fingrom is attributed
to exactly one category (or several if the bill is split, \cref{img:split-bill}), which makes
it possible to answer the following questions at any time: How much did I spend on it? and
How much is left?

\img{img/ch01-s01-budget__split-bill}{Split Bill}{img:split-bill}

\noindent The first category is created during the initial setup wizard. After that, as many
categories as needed can be added from the \textbf{Budgets} page via the \textbf{Add New
Budget Category} button.

\newpage
\subsubsection{Creating a Category}

The form for a budget category contains the following fields (\cref{img:fixed-income}):

\begin{description}
  \item[Title] -- a unique and meaningful name. The \q{/} symbol can be used to group
  categories on the main page: for example, \q{Food/Groceries} and \q{Food/Dining Out} are
  shown together under \q{Food} (\cref{img:main-page}).
  \item[Budgeting Interval] -- the period after which the spent amount is reset:
  \textbf{weekly}, \textbf{monthly} (by default), or \textbf{yearly}.
  \item[Icon] and \textbf{Color} -- visual identification of the category in lists and charts.
  \item[Limit per Month] -- the amount that may be spent within the interval; its meaning
  depends on the value entered (see \cref{budget:limit}).
  \item[Relative Coefficient per Month] -- optional adjustment of the limit for each month of
  the year (see \cref{budget:coef}).
  \item[Currency] -- the currency in which the limit is defined. Bills in other currencies
  are converted automatically using the exchange rates.
  \item[Exclude from Totals] -- when checked, the category is still tracked, but is not
  included in the summary on the main page. This is useful for reimbursable expenses
  (business trips, shared purchases) that should not distort the personal budget.
\end{description}

\img{img/ch01-s01-budget__main-page}{Main Page}{img:main-page}


\subsubsection{Limit Types}
\label{budget:limit}

The value entered in the \textbf{Limit per Month} field (\cref{img:fixed-income}) defines the
category type. The application then highlights the matching type below the input field:

\begin{center}
\begin{tabularx}{\textwidth}{|l|l|X|}
  \hline
  \textbf{Type} & \textbf{Value} & \textbf{Behavior} \\
  \hline
  Unlimited & \q{0} & No limit is set; only the spent amount is recorded and shown. \\
  \hline
  Relative & \q{0 ... 1} & A share of the income received within the current interval (e.g. \q{0.3} means 30\%). \\
  \hline
  Fixed & \q{>= 1} & A fixed amount in the currency of the category. \\
  \hline
\end{tabularx}
\end{center}

\noindent \textbf{Unlimited} categories are useful for expenses that are difficult to plan
for, but which are still worth monitoring. For example, this could be useful during the first
month of using the application, when your actual spending habits are not yet known.

\textbf{Fixed} limits are the most straightforward, with a limit of 400 for groceries meaning
that 400 units of the selected currency can be spent within the interval.

\textbf{Relative} limits are calculated dynamically from all income received since the
beginning of the current interval. Transfers between own accounts are not counted. The income
is then converted into the currency of the category:

\[
  \textit{Limit} = \textit{Coefficient} \times \sum \textit{Incomes}
\]

\noindent For example, if the limit for 'Entertainment' is 0.2 and there are two incomes in
the current month of \$2,000 and \$500, the available amount is $0.2 \times 2\,500 = 500$. If
another income is received later in the month, the limit increases accordingly. Bear in mind
that at the start of the period, before any income has been recorded, the limit is zero.


\subsubsection{Monthly Coefficients}
\label{budget:coef}

Some expenses are seasonal, such as heating in winter, travel in summer and buying gifts in
December. Rather than editing the limit each month, a coefficient can be defined for each
month of the year. To do this, click the \textbf{Declare Limits per Month} button next to the
\textbf{Limit per Month} title (\cref{img:limits-month}). Twelve sliders will appear, one for
each month, each with a value ranging from 0.0 to 4.0. (1.0 by default). The actual limit for
the current month is the base limit multiplied by that month's coefficient:

\img{img/ch01-s01-budget__limits-per-month}{Limits per Month}{img:limits-month}

\begin{center}
\begin{tabular}{|l|r|r|r|}
  \hline
  \textbf{Month} & \textbf{Base Limit} & \textbf{Coefficient} & \textbf{Actual Limit} \\
  \hline
  January  & 100 & 2.00 & 200 \\
  April    & 100 & 1.00 & 100 \\
  July     & 100 & 0.50 &  50 \\
  December & 100 & 1.50 & 150 \\
  \hline
\end{tabular}
\end{center}

\noindent If the coefficient for the current month is greater than one, this is shown next to
the limit in the category description (e.g. '(1.60 coef.)', \cref{img:limits-coef}). To
remove all coefficients, press the \textbf{Cancel Split per Month} button.

\img{img/ch01-s01-budget__limits-coef}{Coefficient per Month}{img:limits-coef}


\subsubsection{Budgeting Intervals}

The interval defines when the spent amount is reset:

\begin{itemize}
  \item \textbf{monthly} -- from the configured \textbf{Start of Month} (any day from 1 to
  31, e.g. the day when the salary arrives);
  \item \textbf{weekly} -- from the configured \textbf{Start of Week} (Monday, Saturday, or
  Sunday);
  \item \textbf{yearly} -- from the 1st of January.
\end{itemize}

\noindent Both start days are defined on the \textbf{Settings} page and affect all categories
simultaneously. Weekly categories are useful for frequent expenses, such as lunches or
groceries, where waiting until the end of the month to realise you have overspent is too
late. Yearly categories are useful for irregular but predictable expenses such as insurance,
vehicle tax and holidays.


\subsection{Working with Categories}
\markboth{Budgeting}{Working with Categories}

\subsubsection{Budgets Page}

The \textbf{Budgets} page lists all active categories. Each row shows the category's title
and a progress bar indicating the amount spent against the limit. For unlimited categories,
the spent amount is shown. Categories with an exceeded limit are highlighted to make
overspending easy to spot.

The header of the page displays the total amount left in all categories, converted into the
default currency. The list can be sorted by total, date of creation or title, and searched
using part of the title (\cref{img:budget-sort}).

\img{img/ch01-s01-budget__sort}{Budget List Sorting}{img:budget-sort}


\subsubsection{Category Details}

Selecting a category opens its details, with a header summarizing the limit, the amount of expenses, and the actual balance, and three tabs (\cref{img:budget-detail}):

\img{img/ch01-s01-budget__budget-detail}{Category Details}{img:budget-detail}

\begin{description}
  \item[Bills] -- all expenses attributed to the category within the current interval, with
  quick access to view or edit each of them.
  \item[Summary] -- spent amounts per month, starting with the current one. Months in which
  the limit was exceeded are marked in red. This is the main source of data for incremental
  budgeting.
  \item[Limit Adjustments] -- the history of changes of the category limit.
\end{description}

\subsubsection{Editing and Deactivating}

A category can be modified at any moment. Changing the limit does not affect the amount
already spent: progress is recalculated against the new limit. For example, if £300 of a £400
limit has been spent (75\%), and the limit is increased to £600, the progress becomes 50\%.

\noindent A category that is no longer needed can be deactivated using the \textbf{Deactivate
Budget Category} option. The category's history is preserved and remains available in
reports, while the category itself is hidden from lists and the selection of a category for
new bills. On touch screens, both actions are also accessible by swiping a row in the list
(\cref{img:swipe}).

\img{img/ch01-s01-budget__swipe}{Swiping a Category Row}{img:swipe}


\subsubsection{Summary Handler}

By default, the total amount shown on the main page and the \textbf{Budgets} page includes
all categories, including those that have been overspent: exceeding the limit reduces the
total. The \textbf{Budget, Summary Handler} option on the \textbf{Settings} page allows you
to change this behavior:

\begin{itemize}
  \item \textbf{Include negative values} -- overspending in one category reduces the overall
  remaining amount, which reflects the actual state of finances;
  \item \textbf{Ignore negative values} -- overspent categories count as zero, so the total
  shows only what is still available in the categories that are within their limits.
\end{itemize}


\subsubsection{Category Prediction}

When a new bill is recorded, Fingrom suggests the most probable budget category. This
suggestion is determined locally on the device using three signals with different weights:
words in the bill description (50\%), the selected account (30\%), and the amount range
(20\%). The more bills that are recorded, the more accurate the prediction becomes
(\cref{img:add-expense}).

\img{img/ch01-s01-budget__add-expense}{Add new Expense}{img:add-expense}


\newpage
\subsection{Monitoring the Budget}
\label{budget:metrics}
\markboth{Budgeting}{Monitoring the Budget}

The \textbf{Metrics} page has a dedicated \textbf{Budgets} tab (\cref{img:budget-metrics}).

\img{img/ch01-s01-budget__budget-metrics}{Budget Metrics}{img:budget-metrics}

\begin{itemize}
  \item \textbf{Forecast Chart} -- the Y-axis shows the budget allocation, where 100\% means
  that all limits have been fully spent, and the X-axis shows the days of the current month.
  The chart consists of two parts: historical data (accumulated expenses so far) and a
  forecast up to the end of the month, calculated using the Monte Carlo simulation method. If
  a forecast line crosses the 100\% mark, this indicates that the limits will be exceeded
  before the end of the month at the current pace.
  \item \textbf{YTD Expenses} -- the expenses accumulated since the beginning of the year,
  which helps to compare the current year with the plan.
  \item \textbf{Table of categories} -- shows the amount spent in each category for the
  current month and the two previous months, making trends easy to spot.
\end{itemize}

\noindent The chart showing how the share of each category changes over time is on the same
page, in the \textbf{Bills} tab. It is called the \textbf{Bar Race for Categories} chart.
You can quickly navigate to the detailed view of a category for a selected month by clicking
on any cell in the table.


\subsection{Practical Tips}
\markboth{Budgeting}{Practical Tips}

\textbf{Start simple.} For the first month, create a few broad, unlimited categories
and record every expense. Only then should you define limits based on the actual numbers.

\textbf{Align the start of the month with your income.} If the salary arrives on the
25th, set 'Start of Month' to 25: the budget period will then match the period in which the
money is actually available.

\noindent \textbf{Prefer relative limits for variable income.} Those with irregular income,
such as freelancers, benefit from limits defined as shares, since they adapt automatically
to good and bad months.

\textbf{Plan for irregular expenses.} Use yearly categories or monthly coefficients
for insurance, taxes, gifts and holidays, so that they do not come as a surprise.

\textbf{Keep a buffer.} An 'Unexpected Expenses' category with a modest limit can
absorb small surprises without disrupting the rest of the plan.

\textbf{Review monthly.} At the end of each period, check the \textbf{Summary} tab
and the \textbf{Metrics} page, and adjust any limits that were consistently exceeded or
never reached.\\
\\

\noindent Common categories to start with:

\begin{multicols}{2}
\begin{itemize}
  \item Housing/Rent
  \item Housing/Utilities
  \item Housing/Maintenance
  \item Food/Groceries
  \item Food/Dining Out
  \item Transport/Fuel
  \item Transport/Public
  \item Health/Insurance
  \item Health/Medical
  \item Leisure/Subscriptions
  \item Leisure/Hobbies
  \item Leisure/Vacations
  \item Personal/Clothing
  \item Education
  \item Unexpected Expenses
\end{itemize}
\end{multicols}
